% Clase 1 — CSP: Propagación de Restricciones y Búsqueda Local
% Lectura: AIMA Cap. 5
\section{CSP Avanzado: Propagación de Restricciones}

\subsection{Consistencia de Arco (AC-3)}

Una variable $X_i$ es \textbf{arco-consistente} con $X_j$ si para cada valor en $D_i$ existe algún valor en $D_j$ que satisface la restricción binaria.

\textbf{AC-3} reduce dominios iterativamente hasta que no quedan inconsistencias o algún dominio queda vacío (sin solución).

\begin{itemize}
    \item Complejidad: $\bigO{cd^3}$ donde $c$ = número de restricciones, $d$ = tamaño del dominio.
\end{itemize}

\subsection{Forward Checking}

Al asignar un valor a $X_i$, eliminar de los dominios de las variables vecinas no asignadas todos los valores inconsistentes.

\subsection{Estructura del Grafo de Restricciones}

Si el grafo de restricciones es un \textbf{árbol}, el CSP puede resolverse en $\bigO{nd^2}$ sin backtracking.

\textbf{Cutset conditioning}: encontrar un conjunto de variables cuya eliminación convierte el grafo en árbol.

\subsection{Búsqueda Local para CSP}

\textbf{Min-conflicts}: en cada paso, reparar la variable con más conflictos eligiendo el valor que minimiza conflictos.

\begin{itemize}
    \item Muy eficiente en la práctica para CSPs grandes.
    \item No garantiza completitud.
\end{itemize}
